\chapter{Coverings}

Now we will study main features of the map $w\:\RR\to \SS^1$,
which was introduced in the previous chapter in order to calculate fundamental group of $\SS^1$.
It will lead us to the definition of covering map and generalization of several statement;
they can be used to caluclate fundamental gropus of other spaces, but more importantly, it will give another connection between algebra and geometry that has applications from both sides.

\section{Definition}

\begin{thm}{Definition}\label{def:evenly-covered}
Let $w\:\tilde{\c{X}}\to \c{X}$ be a continuous map.

An open set $V\subset \c{X}$ is called \index{evenly covered open set}\emph{evenly covered} by $w$ if
\[w^{-1}(V)=\bigcup_{\alpha\in \c{I}}V_\alpha,\]
where $V_\alpha$ are pairwise disjoint open sets in $\tilde{\c{X}}$ such that restriction
$w|_{V_\alpha}$ is a is a homeomorphism $V_\alpha\to V$ for each $\alpha\in \c{I}$.
The sets $V_\alpha$ are called \index{sheet}\emph{sheets} over $V$.

If for every point $x\in \c{X}$ there is an open neighborhood $V\ni x$ that is evenly by $w$, then $w$ is a \index{covering map}\emph{covering map}, $\tilde{\c{X}}$ is a \index{covering space}\emph{covering space} of $\c{X}$.
\end{thm}

Note that $w\:\tilde{\c{X}}\to \c{X}$ is a covering if $\c{X}$ admits an open cover by evenly covered sets.
In particular, from the proof of \ref{prop:lifting}, $w(x)=\exp(2\cdot\pi\cdot i\cdot x)$  defines a covering map $\RR\to\SS^1$ --- this maps and its properties will guide us to study general covering maps.

Most interesting examples of covering maps have path-connected source $\tilde{\c{X}}$ and target  $\c{X}$.
Many authors assume that coverings must be surjective and have connected sourse space.
I might silently assume that this, but formally speaking it is not required (in this text).
For example, one may take several disjoint copies of space $\c{X}$ as $\tilde{\c{X}}$ and map a point in one copy to the corresponding point in $\c{X}$ --- this is a bouring example of a covering map.
More boring examples will appear if $\c{X}$ is not connected; then one may take different number of copies of its connected components and apply the same map.

The following observation says that a covering map must be open.
Note that restriction $w$ to a bounded open interval provides an example of open map which is not covering; so the converse does not hold in general.


\begin{thm}{Observation}\label{thm:covering-open}
Any covering map is an open map.
\end{thm}

\parit{Proof.}
Let $w\:\tilde{\c{X}}\to \c{X}$ be a covering map.
Choose an open set $W\subset \tilde{\c{X}}$ and a point $\tilde x\in W$.

Choose an evenly covered open set $V\ni x=w(\tilde x)$; let $\{V_\alpha\}$ is a family of its sheets.
Since $\tilde x\in w^{-1}(V)$, we have $\tilde x\in V_\alpha$ for some $\alpha$.
Let $V'_\alpha=V_\alpha\cap W$.
Since $w|_{V_\alpha}$ is a homeomorphism, $V'=w(V_\alpha')$ is open in $V$, and (since $V$ is open) $V'$ is open in $\c{X}$.
That is, any $x\in w(W)$ admits an open neighborhood $V'\subset w(W)$;
in other words, $w(W)$ is open.
\qeds

\begin{thm}{Exercise}
Let $w\:\tilde{\c{X}}\to\c{X}$ be a covering map.
Show that the image $w(\tilde{\c{X}})$ is clopen set in $\c{X}$.
In particular, if $\c{X}$ is connected, and $\tilde{\c{X}}\ne\emptyset$, then $w(\tilde{\c{X}})=\c{X}$.
\end{thm}

Let $n$ be an integer.
A covering map $w\:\tilde{\c{X}}\to \c{X}$ is called \emph{$n$-folded} if inverse image of every point $x\in \c{X}$ has exactly $n$ points.

\begin{thm}{Exercise}
Let $w\:\tilde{\c{X}}\to \c{X}$ is a covering map, and let $n$ be an integer.
Assume that $\c{X}$ is connected, and an inverse image of some point $x\in \c{X}$ has exactly $n$ points.
Show that $w$ is $n$-folded.
\end{thm}

\begin{thm}{Exercise}
The following picture shows figure eight and its covering spaces.
Redraw them and mark the arcs of the covering by arrows that correspond to the image in the figure eight.
(The solution is not unique.)
\end{thm}


\begin{figure}[!ht]
\centering
\includegraphics{mppics/pic-1620}
\end{figure}

\section{Real projective plane}

An action $G\acts \c{X}$ is called \index{free action}\emph{free} if for any $x\in\c{X}$ and any $g\in G$, we have
\[g\cdot x=x\quad\iff\quad g=1.\]

\begin{thm}{Proposition}
Let $G\acts \c{X}$ be a continuous free action, let $G$ be finite, and let $\c{X}$ be Hausdorff.
Then the quotient map $\c{X}\to \c{X}/G$ is covering.
\end{thm}

\parit{Proof}
Choose $x\in \c{X}$.
Since $\c{X}$ is Hausdorff we can choose open disjoint sets $V_g\ni x$ and $W_g\ni g\cdot x$ for any $g\ne 1$.
We can assume that $W_g=g\cdot V_g$; if not pass to the intersections $V_g\cap g^{-1}\cdot W_g\ni x$ and
$g\cdot (V_g\cap g^{-1}\cdot W_g)=g\cdot V_g\cap W_g\ni g\cdot x$.
Now consider set
\[V=\bigcap_{g\ne 1}V_g\ni x.\]
Note that $g\cdot V$ are disjoint for all $g\in G$.
The restriction of the quotient map $f$ to $g\cdot V$ is injective, and by \ref{ex:quotient-map:open} it is open.
In other words, $f|_{g\cdot V}$ is a homeomorphism for any $g\in G$, which means that $f$ is a covering.
\qeds

We will denote by $\ZZ_n$ integers modulo $n$;
this is an example of cyclic group of order $n$.
(Recall that all cyclic groups of order $n$ are isomorphic.)

Consider $\ZZ_2$ action on $\SS^2$ generated by the involution $x\mapsto -x$.
The quotient space is called \emph{real protective plane} and denoted by $\RP^2$.
According to the proposition, this gives another intersrting example of covering maps ---  the quotient map $\SS^2\to\RP^2$.
Soon we will use it to calculate $\pi_1(\RP^2)$.




\section{Path liftings}

Recall that in Section~\ref{sec:lifting-s1}, we defined liftings with respect to the covering $\RR\to\SS^1$.
Now we repeat the same definition for general continuous map.

Let $f\:\tilde{\c{X}}\to \c{X}$ and $g\:\c{Y}\to \c{X}$ be continuous maps.
A continuous map $\tilde g\:\c{Y}\to \tilde{\c{X}}$ is called a \index{lift}\emph{lift} of $g$ (with respect to $f$) if
\[f\circ \tilde g=g.\]
In this case we also say that $g$ \index{lifting}\emph{lifts} to $\tilde{\c{X}}$.
In particular, if $\c{Y}=[0,1]$, then $g$ is a path, and the path $\tilde g\:[0,1]\to \tilde{\c{X}}$ is a \emph{lift} of $g$.

The following lemma is a direct generalization of \ref{prop:lifting}, and it can be proved along the same lines.

\begin{thm}{Lemma}\label{lem:path-lifting}
Let $w\:\tilde{\c{X}}\to \c{X}$ be a covering map, and let $a_0\:[0,1]\to \c{X}$ be a path.
Suppose that $\tilde x_0\in \tilde{\c{X}}$ and $x_0\in \c{X}$ be a points such that $w(\tilde x_0)=x_0=a_0(0)$.
Then there is unique lifting $\tilde a_0$ of $a$ such that $\tilde a_0(0)=\tilde x_0$.

Moreover, given a homotopy of paths $a_\tau$ from $a_0$ to $a_1$ there is unique lifting $\tilde a_\tau$ such that $\tilde a_0(0)=\tilde x_0$.
In particular, if $a_0\sim a_1$ then $\tilde a_0\sim \tilde a_1$.
\end{thm}

\begin{thm}{Exercise}\label{ex:induced-hom-covering}
Let $w\:\tilde{\c{X}}\to \c{X}$ be a covering map.
Suppose that $\tilde x_0\in \tilde{\c{X}}$ and $x_0\in \c{X}$ be a points such that $w(\tilde x_0)=x_0$.

\begin{subthm}{ex:induced-hom-covering:injective}
Show that $w_*\:\pi_1(\tilde{\c{X}},\tilde x_0)\to \pi_1(\c{X}, x_0)$ is a \index{monomorphism}\emph{monomorphism} (which means injective homeomorphism).
\end{subthm}

\begin{subthm}{ex:induced-hom-covering:lift}
Let $b$ be a loop in $\c{X}$ with base point $x_0$ and $\tilde b$ be its lift in $\tilde{\c{X}}$ such that $\tilde b(0)=\tilde x_0$.
Show that $\tilde b$ is a loop if and only in $[b]\in \Im w_*$.
\end{subthm}


\begin{subthm}{ex:induced-hom-covering:sc}
Assume  $\tilde{\c{X}}$ is simply connected.
Let $b$ be a loop in $\c{X}$ based at $x_0$ and $\tilde b$  be its lift in $\tilde{\c{X}}$ such that $\tilde b(0)=\tilde x_0$.
Show that $[b]=1\in\pi_1(\c{X},x_0)$ if and only if $\tilde b$ is a loop.
\end{subthm}



\begin{subthm}{ex:induced-hom-covering:endpoint}
Let $a_0$ and $a_1$ be two paths in $\c{X}$ such that $a_0(0)=a_1(0)=x_0$ and $a_0(1)=a_1(1)$.
Suppose $\tilde a_0$ and $\tilde a_1$ be lifts of $a_0$ and $a_1$ such that $\tilde a_0(0)=\tilde a_1(0)=\tilde x_0$.
Show that $\tilde a_0(1)=\tilde a_1(1)$ if and only if $[a_0*\bar a_1]\in w_*(\pi_1(\tilde{\c{X}})$.
\end{subthm}

\end{thm}

\section{Map liftings}

\begin{thm}{Theorem}\label{thm:lifting-criterion}
Let $w\:\tilde{\c{X}}\z\to \c{X}$ be a covering map and let $f\:\c{Y}\z\to \c{X}$ be a continuous map.
Suppose $\c{Y}$ is a path-connected and locally path-connected and $y_0\in \c{Y}$, $x_0\in \c{X}$ and $\tilde x_0\in \tilde{\c{X}}$ be points
such that $w(\tilde x_0)=x_0=f(y_0)$.
Then there exists a unique lift $\tilde f\:\c{Y}\to \tilde{\c{X}}$ such that
$\tilde g(y_0)=\tilde x_0$
if and only if $\Im f_*<\Im w_*$.
\end{thm}

\parit{Proof.}
Assume the lift $\tilde f$ exists, so $f=w\circ \tilde f$.
Then for the induced homeomorphisms, we have $f_*=w_*\circ \tilde f_*$, and hence $\Im f_*<\Im w_*$.



Now choose $y\in \c{Y}$ and a path $a_0$ in $\c{Y}$ from $y_0$ to $y$.
Consider the paths $b_0=f\circ a_0$ in $\c{X}$, and let $\tilde b_0$ be its unique lift such that $\tilde b_0(0)=\tilde x_0$, which is provided by \ref{lem:path-lifting}.

We want to define $\tilde f(y)$ as $\tilde b_0(1)$, but in order to do this we have to check that the point $\tilde b_0(1)\in \tilde{\c{X}}$ does not depend on the choice of path $a_0$ from $y_0$ to $y$.
In other words, if we construct in the same way paths $b_1$ and $\tilde b_1$ for another path $a_1$ from $y_0$ to $y$, then $\tilde b_1(1)=\tilde b_0(1)$.
%Note that $[b_0*\bar b_1]=f_*[a_0*\bar a_1]$. Note that $\widetilde{b_0*\bar b_1}=\tilde b_0* \tilde{\bar b}_1$, where $\tilde{\bar b}_1$ is the lift of $\bar b_1$ such that $\tilde{\bar b}_1(0)= \tilde b_0(1)$. The time-reversed path  $\bar{\tilde{\bar b}}_1$ is another lift of $b_1$. By uniqueness of lift we have that $\bar{\tilde{\bar b}}_1(t)=\tilde b_1(t)$ for any $t$ or $\bar{\tilde{\bar b}}_1(t)\ne\tilde b_1(t)$ for any $t$.
Since $\Im f_*<\Im w_*$, we have that $[b_0*\bar b_1]=f_*[a_0*\bar a_1]$
By \ref{ex:induced-hom-covering:endpoint}, this is equivalent to $[b_0*\bar b_1]\in \Im w_*$, which is the case since $[b_0*\bar b_1]=f_*[a_0*\bar a_1]$ and  $\Im f_*<\Im w_*$.

Now we can define $\tilde f$ by setting $\tilde f(y)\df \tilde b_0(1)$.
By construction $f=w\circ \tilde f$, and it remains to check that $\tilde f$ is continuous.

Fix $y\in \c{Y}$;
let $x\df f(y)$ and $\tilde x=\tilde f(y)$.
Choose an evenly covered open set $V\ni x$.
Let $V_\alpha$ be the unique sheet over $V$ that contains $\tilde x$.

Since $f$ is continuous, we can choose an open neighborhood $W\ni y$ such that $f(W)\subset V$.
Since $\c{Y}$ is locally path-connected, we can assume that $W$ is path-connected.
Choose a path $a_0$ from $y_0$ to $y$
Choose a point $y'\in W$.
Choose a path $a_0$ from $y_0$ to $y$ and a path $a_1$ in $W$ from $y$ to $y'$.
Let $b_0=f\circ a_0$ and $b_1=f\circ a_1$, then $\tilde f(y')=\widetilde{b_0*b_1}(1)$.
Observe that...
\qeds


\section{Lifting lemmas}

Statement of path-lifting lemma (without proof).

Statement about lifting theorem for locally path-connected spaces.

\begin{thm}{Borsuk--Ulam theorem}\label{thm:Borsuk-Ulam}
For every continuous map $f\: \SS^2\to \RR^2$ there exists a pair of antipodal
points $x$ and $-x$ in $\SS^2$ that map to a single point; that is  $f(x)=f(-x)$.
\end{thm}

\parit{Proof.}
Suppose $f(x)\ne f(-x)$ for any $x\in\SS^2$.
Consider the map
\[g(x)=\frac{f(x)- f(-x)}{|f(x)- f(-x)|}.\]
Note that $g$ is a continuous \index{odd map}\emph{odd} map $\SS^2\to\SS^1$; that is $g(-x)=-g(x)$ for any $x\in \SS^2$.

Consider the minimal equivalence relation on $\SS^2$ and $\SS^1$ such that $x\sim -x$.
Let $\RP^2=\SS^2/\sim$

Choose $x\in g^{-1}(1)$, let $f$ be a path from $x$ to $-x$ in $\SS^2$.
Let $w\: \SS^2 \to \SS^2/sim=\RP^2$ is a caovering map.
The compsition $w\circ f$ is a loop in $\RP^2$.


Note that $g\circ f$ is a path from $1$ to $-1$ in $\SS^1$;
its lifting $\widetilde{g\circ f}\: [0,1]\to \RR$ connects a $0$ to a half-integer $n+\tfrac12$ for $n\in \ZZ$.
\qeds

\begin{thm}{Exercse}\label{ex:ABC}
Let $A$, $B$, and $C$ be closed subsets $\SS^2$ such that $A\cup B\cup C=\SS^2$.
Apply the Borsuk--Ulam theorem to the map $f\: \SS^2\to \RR^2$ defined by
\[f\:x\mapsto (\dist_Ax,\dist_Bx)\]
to conclude that one of the sets $A$, $B$, or $C$
contains a pair of antipodal
points in $\SS^2$.

Give an example showing that analogous statement does not hold for four closed set $A$, $B$, $C$, and $D$.

\end{thm}
